\section{Analytic facts used in the limiting argument}
\label{app:analytic-tools}

We record the approximation, localization, and measure arguments
used in Section~\ref{sec:extension}. All gradient norms are
Euclidean. The Sobolev space is inhomogeneous: the function and
its distributional first derivatives belong to $L^1(\R^2)$.
General Sobolev calculus is treated in \cite{EvansGariepy2015};
the exact chain, product, and Lipschitz facts used here also appear
in \cite[Chapter~3, Propositions~3.21--3.22 and
Theorems~3.38--3.39]{Hunter2014}. We give the needed calculations,
including the variation formula that preserves the Euclidean norm.

\subsection{Approximation, translations, and equality sets}

\begin{lemma}\label{lem:appendix-sobolev-calculus}
For real $u\in W^{1,1}(\R^2)$ there are real
$u_j\in C_c^\infty(\R^2)$ converging to $u$ in $W^{1,1}$, and
\begin{equation}\label{eq:appendix-translation}
 \|u(\cdot+z)-u\|_1\le |z|\,\|\nabla u\|_1
 \qquad(z\in\R^2).
\end{equation}
Moreover $|u|\in W^{1,1}(\R^2)$ with
$|\nabla|u||\le|\nabla u|$ almost everywhere. If $\chi$ is
smooth and compactly supported, then $\chi u\in W^{1,1}$ and
$\nabla(\chi u)=\chi\nabla u+u\nabla\chi$.
For locally Lipschitz real functions $v,w$, their maximum is
locally Lipschitz and satisfies, almost everywhere,
\begin{equation}\label{eq:appendix-maximum}
 \nabla\max(v,w)=
 \begin{cases}
  \nabla v,&v>w,\\
  \nabla w,&w>v,\\
  \nabla v=\nabla w,&v=w.
 \end{cases}
\end{equation}
Products of locally Lipschitz functions satisfy the weak product
rule. Their weak derivatives agree almost everywhere with their
classical derivatives.
\end{lemma}

\begin{proof}
The smooth-cutoff product rule follows directly by testing the
distributional derivative of $u$ against $\chi\varphi$.
Choose $\chi_R$ equal to one on $B(0,R)$, zero outside
$B(0,2R)$, with $0\le\chi_R\le1$ and
$|\nabla\chi_R|\le C/R$. Then
\[
 \|\chi_Ru-u\|_1\longrightarrow0,
 \qquad
 \|\nabla(\chi_Ru)-\nabla u\|_1
 \le\int_{|x|>R}|\nabla u|+\frac C R\|u\|_1
 \longrightarrow0.
\]
Convolving $\chi_Ru$ with a smooth compactly supported approximate
identity commutes with its distributional derivatives. The
convolutions converge in $L^1$, for both the function and its
gradient. This last fact follows from continuity of translations
in $L^1$ and the approximate-identity formula; translation
continuity follows first for continuous compactly supported
functions by uniform continuity, and then for $L^1$ functions
by approximation. A diagonal choice of $R$ and convolution
radius proves the asserted $C_c^\infty$ density.

For each smooth approximant, the fundamental theorem of calculus
along a segment gives
\[
 |u_j(x+z)-u_j(x)|
 \le |z|\int_0^1|\nabla u_j(x+sz)|\,\dd s.
\]
Integrating in $x$ and using translation invariance proves the
translation bound for $u_j$. Convergence in $W^{1,1}$ proves
\eqref{eq:appendix-translation} for $u$.

Here is also a direct smooth-chain argument for the absolute
value. For a continuously differentiable function $F$ with bounded
derivative and $F(0)=0$, smooth approximation gives
\[
 \nabla F(u)=F'(u)\nabla u.
\]
Indeed $F(u_j)\to F(u)$ in $L^1$. Along a subsequence $u_j\to u$
almost everywhere, and
\[
 F'(u_j)\nabla u_j-F'(u)\nabla u
 =F'(u_j)(\nabla u_j-\nabla u)
       +(F'(u_j)-F'(u))\nabla u
\]
tends to zero in $L^1$, by boundedness and dominated convergence.
Passing the smooth identity to distributions proves the formula.
Apply it to
$F_\varepsilon(s)=\sqrt{s^2+\varepsilon^2}-\varepsilon$.
These functions satisfy $0\le F_\varepsilon(u)\le|u|$ and
converge to $|u|$ in $L^1$. Their derivatives converge pointwise
to $\operatorname{sgn}(u)$, with value zero at $u=0$, and are
bounded by one. Passing to the limit gives
$\nabla|u|=\operatorname{sgn}(u)\nabla u$ and the claimed norm
inequality.

A locally Lipschitz function restricts to a Lipschitz, hence
absolutely continuous, function on each coordinate line in a
rectangle. One-dimensional integration by parts and Fubini's
theorem give its weak partial derivatives. Rademacher's theorem
identifies these with its classical gradient almost everywhere
\cite[Chapter~3, Theorems~3.38--3.39]{Hunter2014}. The usual
product rule at joint differentiability points, or equivalently
the one-dimensional rule on those lines, proves the weak
Lipschitz product rule.

For the equality set in \eqref{eq:appendix-maximum}, put $a=v-w$.
At almost every point of $\{a=0\}$ the set has density one and
$a$ is differentiable. Its derivative there is zero: a nonzero
derivative would give a cone of fixed positive relative area on
which $a$ is nonzero in every sufficiently small ball, contrary
to density one. Thus $\nabla v=\nabla w$ almost everywhere on
their equality set. At such a point,
$a(x+h)=o(|h|)$ implies $a_+(x+h)=o(|h|)$, so $a_+$ also has
derivative zero. On the open sets $a>0$ and $a<0$, its gradient
is respectively $\nabla a$ and zero. Since
$\max(v,w)=w+a_+$, this proves \eqref{eq:appendix-maximum}.
\end{proof}

\begin{lemma}[Planar Sobolev embedding]\label{lem:appendix-planar-embedding}
For every real $u\in W^{1,1}(\R^2)$,
\[
 \|u\|_2^2\le\|\partial_1u\|_1\|\partial_2u\|_1,
 \qquad\text{in particular }\|u\|_2\le\|\nabla u\|_1.
\]
\end{lemma}

\begin{proof}
For $u\in C_c^\infty$, set
\[
 a(y)=\int_\R|\partial_1u(x,y)|\,\dd x,
 \qquad b(x)=\int_\R|\partial_2u(x,y)|\,\dd y.
\]
The fundamental theorem of calculus on the two coordinate lines
gives $|u(x,y)|\le a(y)$ and $|u(x,y)|\le b(x)$.
Thus $|u(x,y)|^2\le a(y)b(x)$; integrating proves the first
inequality. For general $u$, approximate as in
Lemma~\ref{lem:appendix-sobolev-calculus}. Applying the smooth
inequality to differences shows that the approximants are Cauchy
in $L^2$. Their $L^2$ limit agrees almost everywhere with their
$L^1$ limit $u$, by passage to a common almost-everywhere
convergent subsequence. Pass to the limit in the inequality.
Each coordinate derivative has $L^1$ norm at most
$\|\nabla u\|_1$, giving the last assertion.
\end{proof}

\subsection{Square localization}

\begin{lemma}\label{lem:appendix-square}
For a square $Q$ of side length $t>0$ and
$u\in W^{1,1}(\R^2)$, write $u_Q=|Q|^{-1}\int_Qu$. Then
\begin{equation}\label{eq:appendix-poincare}
 \int_Q|u-u_Q|\le Ct\int_Q|\nabla u|.
\end{equation}
For every grid square $Q$ of side $t$, there is a smooth cutoff
$\chi_Q$ equal to one on the closed concentric dilation $3Q$,
supported in the interior of $5Q$, with $0\le\chi_Q\le1$ and
$|\nabla\chi_Q|\le C/t$. With $c_Q=u_{5Q}$, it satisfies
\begin{equation}\label{eq:appendix-localization}
 \|\nabla[\chi_Q(u-c_Q)]\|_1
       \le C\int_{5Q}|\nabla u|.
\end{equation}
Every disk of radius at most $t$ centered in $Q$ lies in $3Q$.
The dilations $5Q$ of grid squares have bounded overlap. If the
squares under consideration meet
$O\subset\{x:\dist(x,K)<t\}$, all their dilations $5Q$ lie in
$\{x:\dist(x,K)<6t\}$. Constants are independent of
$Q,t,u,O,K$.
\end{lemma}

\begin{proof}
By rotation and translation it suffices to take $Q=(0,t)^2$.
First suppose $u$ is smooth. Jensen's inequality gives
\[
 \int_Q|u-u_Q|
 \le t^{-2}\int_Q\int_Q|u(x)-u(y)|\,\dd y\,\dd x.
\]
Join $x=(x_1,x_2)$ to $(y_1,x_2)$ and then to $y=(y_1,y_2)$.
For the first segment, with $x_2$ fixed, integrating its
derivative bound in $x_1,y_1$ gives at most
\[
 t^2\int_0^t|\partial_1u(s,x_2)|\,\dd s.
\]
Indeed each $s$ lies between $x_1$ and $y_1$ for a set of pairs
of area at most $t^2$. Integration in $x_2$ and $y_2$ therefore
gives $t^3\int_Q|\partial_1u|$. The second segment similarly
costs at most $t^3\int_Q|\partial_2u|$.
After division by $t^2$, the bound
$|\partial_1u|+|\partial_2u|\le\sqrt2|\nabla u|$ proves
\eqref{eq:appendix-poincare}. Global smooth approximation proves
it for the stated Sobolev input, since both the functions and
their means converge on $Q$.

Rescale a fixed smooth cutoff to obtain $\chi_Q$. The product
rule and \eqref{eq:appendix-poincare}, now on $5Q$, yield
\[
 \|\nabla[\chi_Q(u-c_Q)]\|_1
 \le\int_{5Q}|\nabla u|+\frac C t\int_{5Q}|u-c_Q|
 \le C\int_{5Q}|\nabla u|.
\]
In each coordinate the disk containment follows from
$t/2+t=3t/2$. A point lies in only a fixed number of the
dilated squares, since their centers form a grid of spacing $t$.
Finally, if $x\in Q\cap O$ and $y\in5Q$, then
$|y-x|\le3\sqrt2t$. Hence
$\dist(y,K)<(1+3\sqrt2)t<6t$, as claimed.
\end{proof}

\subsection{Vector variation and local lower semicontinuity}

For $u\in L^1_{\mathrm{loc}}(\R^2)$ and open $O\subset\R^2$,
define its distributional variation on $O$ by
\begin{equation}\label{eq:appendix-variation}
 \mathcal V(u;O)=
 \sup\left\{-\int u\,\operatorname{div}\varphi:
       \varphi\in C_c^1(O;\R^2),\ |\varphi(x)|\le1\right\}.
\end{equation}
The test-field norm is Euclidean. The class is closed under
negation, so the same supremum results if the tested integral
is replaced by its absolute value.

\begin{lemma}\label{lem:appendix-vector-measure}
Suppose $u_j\in W^{1,1}_{\mathrm{loc}}(\R^2)$ converge locally
in $L^1$ to $u$, and
$\sup_j\int_{\R^2}|\nabla u_j|\le M<\infty$.
Then the distributional gradient of $u$ is a finite vector Radon
measure $H$, with
\begin{equation}\label{eq:appendix-local-lsc}
 |H|(\R^2)\le M,\qquad
 |H|(O)=\mathcal V(u;O)
       \le\liminf_j\int_O|\nabla u_j|
 \quad\text{for every open }O.
\end{equation}
If $|H|(K)=0$ for every compact Lebesgue-null set $K$, then
$H=G\,\dd x$ for some $G\in L^1(\R^2;\R^2)$ and
$\|G\|_1=|H|(\R^2)$. In that case
$u\in W^{1,1}_{\mathrm{loc}}$, with weak gradient $G$.
\end{lemma}

\begin{proof}
For every smooth compactly supported vector field $\varphi$,
local $L^1$ convergence implies
\[
 \left|-\int u\,\operatorname{div}\varphi\right|
 =\lim_j\left|\int\nabla u_j\cdot\varphi\right|
 \le M\|\varphi\|_\infty.
\]
Smooth compactly supported fields are uniformly dense in
$C_0(\R^2;\R^2)$, the continuous fields vanishing at infinity:
first cut off outside a large compact set and then convolve.
Thus the functional extends uniquely to this space with norm
at most $M$. Apply the scalar Riesz representation theorem to
its two coordinates to obtain finite signed Radon measures
$H_1,H_2$. Their vector $H=(H_1,H_2)$ represents the
functional and therefore the distributional gradient of $u$.
The Riesz and Radon--Nikodym theorems used here are recorded in
\cite[Chapter~1, Theorem~5.14 and Equation~(4.17)]{Simon2018}.

We spell out why the representation retains the Euclidean total
variation rather than the sum of coordinate variations. Put
$\lambda=|H_1|+|H_2|$ and write $H=q\lambda$ using the scalar
Radon--Nikodym theorem. The coordinate densities of $q$ have
absolute value at most one. For a Borel set $E$, the definition
of vector total variation by finite Borel partitions gives
\[
 |H|(E)=\int_E|q|\,\dd\lambda.
\]
The upper bound follows from the triangle inequality on each
part. For the reverse bound, partition the bounded range of $q$
into finitely many sets of diameter at most $\varepsilon$, and
choose a vector $q_s$ in each. On the corresponding partition
$E_s$ of $E$,
\[
 \sum_s|H(E_s)|
 \ge\sum_s|q_s|\lambda(E_s)-\varepsilon\lambda(E)
 \ge\int_E|q|\,\dd\lambda-2\varepsilon\lambda(E).
\]
Let $\varepsilon\downarrow0$. Set $\mu=|H|$ and
$\sigma=q/|q|$ where $q\ne0$, with value zero elsewhere;
then $H=\sigma\mu$ and $|\sigma|=1$ almost everywhere for $\mu$.

For every open $O$, this representation gives
\begin{equation}\label{eq:appendix-continuous-variation}
 \mu(O)=\sup\left\{\int\varphi\cdot\dd H:
          \varphi\in C_c(O;\R^2),\ |\varphi|\le1\right\}.
\end{equation}
One inequality is immediate. For the other, approximate $\sigma$
on $O$ in $L^1(\mu)$ by simple vector functions with values in
the closed unit ball. Inner regularity permits replacement of
their finitely many disjoint measurable sets by disjoint compact
subsets of $O$, losing arbitrarily little $\mu$-mass. Choose
continuous cutoff functions equal to one on these compact sets,
with mutually disjoint compact supports in $O$, and multiply by
the corresponding vector values. Their sum has norm at most
one and approximates $\sigma$ in $L^1(\mu)$, with arbitrarily
small error. This proves the reverse inequality. These are the
usual polar and open-set variation formulas; see also
\cite[Chapter~2, Equations~(2.1)--(2.4)]{Simon2018}.

Each continuous compactly supported test field in
\eqref{eq:appendix-continuous-variation} can be extended by zero
and convolved with a nonnegative smooth probability kernel of
sufficiently small support. The resulting smooth fields remain
supported in $O$, have norm at most one, and converge uniformly.
Consequently the same supremum is obtained with $C_c^1$ test
fields. The distributional identity identifies it with
$\mathcal V(u;O)$ in \eqref{eq:appendix-variation}. Taking
$O=\R^2$ and using the functional norm bound gives
$|H|(\R^2)\le M$.

For each fixed admissible field supported in $O$,
\[
 -\int u\,\operatorname{div}\varphi
 =\lim_j\int\nabla u_j\cdot\varphi
 \le\liminf_j\int_O|\nabla u_j|.
\]
Taking the supremum proves the local lower semicontinuity in
\eqref{eq:appendix-local-lsc}. This argument requires only local
integrability of $u$.

Finally, if $\mu(K)=0$ for all compact Lebesgue-null $K$, inner
regularity implies $\mu(E)=0$ for every Borel null set $E$.
Every Lebesgue-null set is contained in a Borel null set, so
$\mu$ is absolutely continuous with respect to Lebesgue measure.
Write $\mu=h\,\dd x$ by the positive Radon--Nikodym theorem,
where $h\ge0$ and $\int h=\mu(\R^2)$. Set $G=h\sigma$,
with value zero where $h=0$. Then $H=G\,\dd x$ and
$|G|=h$ almost everywhere. Thus
$\|G\|_1=|H|(\R^2)$, and the defining distributional identity
for $H$ is precisely the weak-gradient identity for $G$.
The needed inner regularity is, for example,
\cite[Chapter~1, Lemma~5.7]{Simon2018}.
\end{proof}

In Section~\ref{sec:extension}, this lemma first yields a finite
derivative measure from the expanding bands. The signed cutoff
argument then proves that it vanishes on compact null sets;
the lemma's final assertion identifies the resulting weak gradient.
